\begin{center}
\centering
\providecommand{\hlspan}[1]{{\setlength{\fboxsep}{1.5pt}\colorbox{ThesisPrimary!24}{\textcolor{ThesisPrimaryDark}{\textbf{#1}}}}}
\begin{tikzpicture}
  \node[draw=ThesisNeutral, fill=ThesisPaper, rounded corners=3pt, line width=0.5pt,
        inner sep=10pt] {%
    \begin{minipage}[t]{0.82\textwidth}
      {\scriptsize\sffamily\color{ThesisNeutral} COMPTE-RENDU ANATOMOPATHOLOGIQUE $\cdot$ fictif}\\[4pt]
      {\color{ThesisNeutral}\rule{\linewidth}{0.3pt}}\\[5pt]
      \small\itshape\color{ThesisInk}%
      \textbf{Nature du pr\'el\`evement.} Biopsies.\\[4pt]
      \textbf{Examen microscopique.} Prolif\'eration tumorale faite d'\'el\'ements
      tass\'es, noyaux hyperchromatiques.\\[4pt]
      \textbf{Immunohistochimie.} Les cellules tumorales expriment l'\hlspan{AE1AE3},
      la \hlspan{chromogranine}, la \hlspan{synaptophysine} et le \hlspan{CD56}.
    \end{minipage}%
  };
\end{tikzpicture}

\smallskip
{\scriptsize\sffamily\color{ThesisNeutral}%
 biomarqueurs $\cdot$ {\raise0.5pt\hbox{$\sim$}}9 mentions par document $\cdot$ registre court}

\captionof{figure}{Un compte-rendu anatomopathologique fictif de PARHAF, avec ses
mentions de biomarqueurs surlign\'ees.}
\label{fig:p3-parhaf-doc}
\end{center}
